\chapter{CONCLUSÃO}

Os resultados obtidos demonstram a eficácia dos métodos de estimativa de massa estelar aplicados a diferentes aglomerados estelares. Os modelos de regressão para a relação massa-magnitude mostraram-se robustos, com alta correlação entre as massas previstas e observadas, validado pelos testes estatísticos. O método de ajuste de isócronas apresentou limitações, principalmente para estrelas de maior massa, com uma tendência à subestimação em certos aglomerados. Por outro lado, o método de relação massa-magnitude, utilizando o modelo BHAC15, forneceu resultados mais precisos consistentemente, especialmente para estrelas de baixa e média massa, embora apresente maior dispersão para objetos mais massivos e para aglomerados que tenham uma dispersão maior em idade.

A aplicação desses métodos aos diferentes aglomerados revelou que, para o ajuste da relação massa-magnitude utilizando modelos teóricos de isócronas, o modelo baseado em K-ésimos Vizinhos Próximos (KNN) foi o que obteve melhores métricas em termos de precisão geral, enquanto o ajuste de isócronas é mais preciso para estimativas em aglomerados mais jovens ou com uma heterogeneidade maior nas idades. Assim, a relação massa-magnitude mostrou-se mais adaptável às condições variáveis dos aglomerados estudados, oferecendo estimativas de massa mais confiáveis para uma ampla gama de massas estelares, precisando de menos informação para ser utilizado.